% SGA 3, Expose VI_B, Loop 1 English reconstruction.
% R3: independent convention-v2 PDF replay assigned distinct
% precise anchors to the two custom-labelled conclusions 11.15(i)/(ii).
% Scope: Lemma 11.12 through Theorem 11.16, including the
% complete proof of Theorem 11.16 and editor's notes 125--129.
% Authority: Polo--Gille Exp6B-13oct24.pdf, component-local
% pages 90--94 (combined-reader pages 431--435; re-edition
% headers 420--424; original-print margin markers 405--409).
% Lemma 11.17.0 begins on local page 94 / combined page 435 /
% header 424 under original-print marker 409 and is excluded.
% OCR is locator-only. Reinhold e7a259f English Markdown is
% comparison-only and is credited in the component report.
% No source diagram occurs in this range.


\medskip
\phantomsection\label{lem:sga3-VIB-11.12}%
\noindent\textbf{引理 11.12.---}
\emph{设 \(k\) 为域，\(G\) 为仿射 \(k\)-群。置 \(A=\mathcal A(G)\)。给定 \(x\in A\)，存在 \(A\) 的一个有限生成 \(k\)-子代数 \(B\)，使得}
\[
  x\in B,\qquad
  \Delta(B)\subset B\otimes_kB,\qquad
  u(B)\subset B,
\]
\emph{其中 \(u\) 表示与 \(G\) 的取逆态射相对应的 \(A\) 的对合。}%
\footnote{\label{note:sga3-VIB-125}编者注：一方面，\hyperref[sec:sga3-VIB-13]{第 13 节}将此推广到 \(G\) 在维数至多为 \(2\) 的正则底上仿射且平坦的情形。另一方面，下面扩充并改正了原文。}

可以假设 \(x\neq0\)。此时 \(\Delta(x)\neq0\)，因为
\[
  (\varepsilon\otimes\mathrm{id})\Delta(x)=x,
\]
其中 \(\varepsilon\) 表示 \(A\) 的增广\textup{（余单位）}。写成
\phantomsection\label{form:sga3-VIB-11.12-eq1}%
\[
  \tag{1}
  \Delta(x)=\sum_{j=1}^{n}e_j\otimes a_j,
  \qquad n\ \text{最小}.
\]
于是各 \(e_j\)\textup{（分别地，各 \(a_j\)）}线性无关。将 \((e_1,\ldots,e_n)\) 补成 \(A\) 的一组基 \((e_i)_{i\in I}\)，并对
\[
  j\in J=\{1,\ldots,n\},
\]
写成
\[
  \Delta(e_j)=\sum_{i\in I}e_i\otimes b_{ij}.
\]
对 \hyperref[form:sga3-VIB-11.12-eq1]{\textup{(1)}} 分别施以 \(\Delta\otimes\mathrm{id}\) 和 \(\mathrm{id}\otimes\Delta\)，由 \hyperref[def:sga3-VIB-11.8.0]{11.8.0} 的公理 \hyperref[ax:sga3-VIB-11.8.0-CO1]{\textup{(CO 1)}}\textup{（亦见第 I 篇 \SGARefLink{sga3:i:subsection:4-2}{4.2} 的 \textup{(HA 1)}）} 得到
\[
  \sum_{j\in J}e_j\otimes\Delta(a_j)
  =
  \sum_{\ell\in J}\Delta(e_\ell)\otimes a_\ell
  =
  \sum_{i\in I}e_i\otimes
    \left(\sum_{\ell\in J}b_{i\ell}\otimes a_\ell\right).
\]
由于各 \(e_i\) 线性无关，可得
\phantomsection\label{form:sga3-VIB-11.12-eq2}%
\[
  \tag{2}
  \Delta(a_j)=\sum_{\ell\in J}b_{j\ell}\otimes a_\ell
  \qquad (j\in J).
\]
设 \(B\) 为由各 \(b_{ij}\) 和 \(u(b_{ij})\)（\(i,j\in J\)）生成的 \(A\) 的有限生成 \(k\)-子代数。显然 \(u(B)=B\)。

对 \hyperref[form:sga3-VIB-11.12-eq2]{\textup{(2)}} 分别施以 \(\Delta\otimes\mathrm{id}\) 和 \(\mathrm{id}\otimes\Delta\)，由 \textup{(CO 1)} 再次得到等式
\[
  \sum_{\ell\in J}\Delta(b_{j\ell})\otimes a_\ell
  =
  \sum_{i\in J}b_{ji}\otimes\Delta(a_i)
  =
  \sum_{i,\ell\in J}b_{ji}\otimes b_{i\ell}\otimes a_\ell.
\]
由于各 \(a_\ell\) 线性无关，得到
\phantomsection\label{form:sga3-VIB-11.12-eq3}%
\[
  \tag{3}
  \Delta(b_{j\ell})
    =\sum_{i\in J}b_{ji}\otimes b_{i\ell}
  \qquad (j,\ell\in J).
\]
此外，
\[
  \Delta\circ u=(u\otimes u)\circ v\circ\Delta,
  \qquad v(a\otimes b)=b\otimes a,
\]
从而
\phantomsection\label{form:sga3-VIB-11.12-eq4}%
\[
  \tag{4}
  \Delta\bigl(u(b_{j\ell})\bigr)
    =\sum_{i\in J}u(b_{i\ell})\otimes u(b_{ji})
  \qquad (j,\ell\in J).
\]
由于 \(\Delta\) 是代数同态，\hyperref[form:sga3-VIB-11.12-eq3]{\textup{(3)}} 和 \hyperref[form:sga3-VIB-11.12-eq4]{\textup{(4)}} 蕴含
\[
  \Delta(B)\subset B\otimes_kB.
\]
最后，将 \hyperref[def:sga3-VIB-11.8.0]{11.8.0} 的公理 \hyperref[ax:sga3-VIB-11.8.0-CO2]{\textup{(CO 2)}}\textup{（亦见第 I 篇 \SGARefLink{sga3:i:subsection:4-2}{4.2} 的 \textup{(HA 2)}）} 应用于 \hyperref[form:sga3-VIB-11.12-eq2]{\textup{(2)}} 和 \hyperref[form:sga3-VIB-11.12-eq1]{\textup{(1)}}，得到%
\begingroup
\renewcommand{\thefootnote}{\fnsymbol{footnote}}%
\footnote[1]{\label{note:sga3-VIB-11.12-counit-source-reading}
源文读法注：法文重版本印作 \(a_j=\sum_{i\in I}\varepsilon(a_j)b_{ij}\)。对 \hyperref[form:sga3-VIB-11.12-eq2]{\textup{(2)}} 施以 \(\mathrm{id}\otimes\varepsilon\)，得到这里列出的第一个恒等式。Reinhold 对照版本沿袭了所印公式。}
\endgroup
\phantomsection\label{form:sga3-VIB-11.12-counit-identities}%
\[
  a_j=\sum_{i\in J}b_{ji}\varepsilon(a_i),
  \qquad
  x=\sum_{j\in J}\varepsilon(e_j)a_j.
\]
于是 \(x\in B\)，引理得证。

\medskip
\phantomsection\label{prop:sga3-VIB-11.13}%
\noindent\textbf{命题 11.13.---}
\emph{设 \(k\) 为域，\(G\) 为代数为 \(A\) 的仿射 \(k\)-群。则 \(G\) 是由有限型仿射 \(k\)-群组成的一个递增滤过系统的射影极限，其过渡态射忠实平坦。}

若 \(B\) 和 \(B'\) 是 \(A\) 的两个在 \(\Delta\) 和 \(u\) 下稳定的有限生成子代数，则由 \(B\) 和 \(B'\) 生成的子代数也具有这些性质。因此，由\hyperref[lem:sga3-VIB-11.12]{引理 11.12}，\(A\) 是一个在 \(\Delta\) 和 \(u\) 下稳定的有限生成子代数的递增滤过族 \((B_i)_{i\in I}\) 的归纳极限。每个 \(B_i\) 配备 \(u\) 的限制以及由 \(\Delta\) 诱导的态射
\[
  B_i\longrightarrow B_i\otimes_kB_i
\]
后，都是 Hopf 代数。因此，由第 I 篇 \SGARefLink{sga3:i:subsection:4-2}{4.2}，它是一个在 \(k\) 上有限型的仿射 \(k\)-群 \(G_i\) 的代数。最后，由于
\[
  A=\varinjlim_i B_i,
\]
有
\[
  G=\varprojlim_i G_i
\]
\textup{（参见 EGA IV\(_3\)，8.2.3）}。由下述引理，过渡态射忠实平坦。

\medskip
\phantomsection\label{lem:sga3-VIB-11.14}%
\noindent\textbf{引理 11.14.---}
\emph{设 \(k\) 为域，\(u:G\to H\) 为有限型仿射 \(k\)-群之间的态射，并设}
\[
  u^\natural:B\longrightarrow A
\]
\emph{为相应的 \(k\)-代数态射。}%
\footnote{\label{note:sga3-VIB-126}编者注：我们用 \(u^\natural\)\textup{（代替 \(u^\circ\)）}表示与 \(u:\operatorname{Spec}(A)\to\operatorname{Spec}(B)\) 相对应的态射 \(B\to A\)。}
\emph{则 \(u\) 忠实平坦，当且仅当 \(u^\natural\) 单射。}

此条件显然必要\textup{（参见 EGA 0\(_{\mathrm I}\)，6.6.1）}。反之，假设 \(u^\natural\) 单射，并置
\[
  N=\operatorname{Ker}(u).
\]
由 \hyperref[thm:sga3-VIA-3.3.2]{VI\(_A\)，3.3.2} 和 \hyperref[prop:sga3-VIA-5.4.1]{VI\(_A\)，5.4.1}，\(G/N\) 是有限型 \(k\)-群，且 \(u\) 分解为
\[
  G\xrightarrow{\ p\ }G/N\xrightarrow{\ v\ }H,
\]
其中 \(p\) 忠实平坦，\(v\) 为闭浸入。由于 \(H\) 仿射，\(G/N\) 也仿射，而且
\[
  v^\natural:B\longrightarrow\mathcal O(G/N)
\]
满射\textup{（参见 EGA I，4.2.3）}。但
\[
  u^\natural=p^\natural\circ v^\natural
\]
单射，故 \(v^\natural\) 也单射。因此 \(v^\natural\) 是同构，从而 \(v\) 也是同构。由于 \(p\) 忠实平坦，\(u\) 也忠实平坦。

\medskip
\phantomsection\label{def:sga3-VIB-11.15}%
\noindent\textbf{定义 11.15.---}%
\footnote{\label{note:sga3-VIB-127}编者注：我们增添了 \(G\) 拟紧的假设，并详细说明了条件 \textup{(i)} 与 \textup{\SGARefLink{concl:sga3-VIB-11.15-i}{(ii)}} 的等价性。}
\emph{设 \(k\) 为域，\(G\) 为拟紧 \(k\)-群，\(V\) 为带有 \(G\) 的 \(k\)-线性作用的 \(k\)-向量空间，因而由 \hyperref[prop:sga3-VIB-11.6.1]{11.6.1} 带有一个 \(\mathcal A(G)\)-余模结构}
\[
  \delta:V\longrightarrow V\otimes\mathcal A(G).
\]
\emph{设 \(v\in V\) 非零。下列条件等价：}
\begin{enumerate}
\item[\textup{(i)}]\phantomsection\label{concl:sga3-VIB-11.15-i}
  \emph{存在一个必然唯一的 \(\lambda\in\mathcal A(G)\)，使得}
\[
    \delta(v)=v\otimes\lambda.
  \]
\item[\textup{(ii)}]\phantomsection\label{concl:sga3-VIB-11.15-ii}
  \emph{对每个 \(k\)-代数 \(R\) 及每个 \(g\in G(R)\)，有 \(g\cdot v\in Rv\)；即存在一个必然唯一的 \(f\in R\)，使得在 \(V\otimes R\) 中，}
\[
    g\cdot v=v\otimes f.
  \]
\end{enumerate}

显然 \hyperref[concl:sga3-VIB-11.15-i]{\textup{(i)}} 蕴含 \hyperref[concl:sga3-VIB-11.15-ii]{\textup{(ii)}}。反之，若 \textup{\SGARefLink{concl:sga3-VIB-11.15-i}{(ii)}} 成立，取
\[
  R=\mathcal A(G),
  \qquad
  g=\operatorname{id}_{\mathcal A(G)}.
\]
应用该条件。此时存在唯一的 \(\lambda\in\mathcal A(G)\)，使 \(\delta(v)=v\otimes\lambda\)。

当 \(v\) 满足这些条件时，称 \(v\) 为\emph{关于 \(G\) 的半不变向量}，\(\lambda\) 为 \(v\) 的\emph{权}。也称 \(v\) 为\emph{权为 \(\lambda\) 的半不变量}。

用 \(\Delta\) 表示 \(\mathcal A(G)\) 的余乘法。等式
\phantomsection\label{form:sga3-VIB-11.15-group-like}%
\[
  v\otimes\lambda\otimes\lambda
  =
  (\delta\otimes\mathrm{id})(\delta(v))
  =
  (\mathrm{id}\otimes\Delta)(\delta(v))
  =
  v\otimes\Delta(\lambda)
\]
蕴含
\[
  \Delta(\lambda)=\lambda\otimes\lambda.
\]
因此，\(\lambda\) 定义一个 Hopf 代数态射
\[
  \mathcal A(\mathbb G_{m,k})
    =k[T,T^{-1}]
    \longrightarrow\mathcal A(G),
  \qquad T\longmapsto\lambda,
\]
从而定义一个 \(k\)-群态射
\[
  \lambda:G\longrightarrow\mathbb G_{m,k}.
\]
所以 \(\lambda\) 是 \(G\) 的一个\emph{特征标}，称为\emph{与半不变向量 \(v\) 相伴的特征标}。

\medskip
\phantomsection\label{lem:sga3-VIB-11.16.0}%
\noindent\textbf{引理 11.16.0.---}%
\footnote{\label{note:sga3-VIB-128}编者注：增补了这一引理，取自 \textup{[DG70]}，第 II.2 节，\SGARefLink{prop:sga3-VIB-3.5}{3.5}。}
\emph{设 \(k\) 为域，\(H\) 为仿射 \(k\)-群，\(V\) 为维数为 \(n\) 的 \(H\)-模，\(U\) 为 \(V\) 的 \(d\) 维向量子空间。考虑直线}
\[
  D=\bigwedge\nolimits^dU
    \subset\bigwedge\nolimits^dV.
\]
\emph{则 \(U\) 在 \(H\) 下稳定，当且仅当 \(D\) 在 \(H\) 下稳定。}

必要性显然。为证明充分性，可假设 \(d<n\)。设 \((e_1,\ldots,e_d)\) 为 \(U\) 的一组基，并将其补成 \(V\) 的一组基 \((e_1,\ldots,e_n)\)。对每个 \(k\)-代数 \(R\)，模
\[
  V_R=V\otimes R
\]
在 \(R\) 上自由，而且
\phantomsection\label{form:sga3-VIB-11.16.0-UR}%
\[
  U_R
  =
  \{\,v\in V_R\mid
    v\wedge(e_1\wedge\cdots\wedge e_d)=0\,\}.
\]
事实上，对 \(i>d\)，元素
\[
  e_i\wedge e_1\wedge\cdots\wedge e_d
\]
在 \(\bigwedge_R^{d+1}V_R\) 中线性无关。由于 \(H(R)\) 按照
\[
  h(x_1\wedge\cdots\wedge x_s)
  =
  h(x_1)\wedge\cdots\wedge h(x_s),
\]
作用于 \(\bigwedge_R^\bullet V_R\)，可知若 \(H(R)\) 稳定
\[
  D_R=R(e_1\wedge\cdots\wedge e_d),
\]
则它也稳定 \(U_R\)。

\medskip
\phantomsection\label{thm:sga3-VIB-11.16}%
\noindent\textbf{定理 11.16 \textup{(Chevalley)}.---}
\emph{设 \(k\) 为域，\(G\) 为仿射代数 \(k\)-群，\(H\) 为 \(G\) 的闭子群概形。}%
\footnote{\label{note:sga3-VIB-129}编者注：首先回顾，\(G\) 的每个子群概形都闭\textup{（\hyperref[cor:sga3-VIB-1.4.2]{1.4.2}）}。其次，我们以通常的形式陈述此结果，即“\(H\) 是 \(G\) 的一个表示中的某条直线的稳定子”，同时保留原文以 \(\mathcal A(G)\) 中的一列半不变量 \(a_1,\ldots,a_n\) 表述的形式。}
\emph{则存在有限维 \(G\)-模 \(V\) 及 \(V\) 中的一条直线 \(D\)，使得}
\[
  H=\operatorname{Norm}_G(D);
\]
\emph{即对每个 \(k\)-代数 \(R\)，}
\phantomsection\label{form:sga3-VIB-11.16-stabilizer}%
\[
  H(R)=\{\,g\in G(R)\mid g(D_R)=D_R\,\}.
\]
\emph{等价地，存在有限多个元素 \(a_1,\ldots,a_n\in\mathcal A(G)\)，它们关于 \(H\) 的“右”作用都是具有同一权 \(\lambda\) 的半不变量；即}
\[
  a_i(gh)=\lambda(h)a_i(g)
\]
\emph{对每个 \(k\)-代数 \(R\)、\(g\in G(R)\) 及 \(h\in H(R)\) 都成立，并且 \(H\) 是使各 \(a_i\) 为半不变量的 \(G\) 的最大闭子群概形。}

\emph{证明。}
用 \(\Delta\) 和 \(\varepsilon\) 分别表示
\[
  A=\mathcal A(G).
\]
的余乘法和增广。存在 \(A\) 的一个包含于 \(\operatorname{Ker}\varepsilon\) 的理想 \(I\)，使得
\[
  H=\operatorname{Spec}(A/I),
  \qquad
  \Delta(I)\subset I\otimes A+A\otimes I.
\]
置 \(B=A/I\)，并设 \(\pi:A\to B\) 为投影。考虑 \(H\) 在 \(A\) 上由
\[
  (h\phi)(g)=\phi(gh).
\]
给出的左作用。相应的 \(B\)-余模结构为
\phantomsection\label{form:sga3-VIB-11.16-H-coaction-A}%
\[
  \overline{\Delta}:\quad
  A\xrightarrow{\ \Delta\ }A\otimes A
    \xrightarrow{\ \mathrm{id}_A\otimes\pi\ }A\otimes B.
\]
此时 \(I\) 是 \(A\) 的 \(H\)-子模，因为
\[
  \overline{\Delta}(I)\subset I\otimes B.
\]
\(k\)-代数 \(A\) 有限生成，因而是 Noether 环，故 \(I\) 有有限生成元组 \((x_1,\ldots,x_r)\)。由 \hyperref[lem:sga3-VIB-11.8]{11.8}，各 \(x_i\) 包含在 \(A\) 的一个有限维 \(G\)-子模 \(V\) 中。于是
\[
  W=V\cap I
\]
是有限维 \(H\)-模；记
\[
  d=\dim_kW.
\]
由于 \(V\) 包含每个 \(x_i\)，\(W\) 生成理想 \(I\)。

置
\[
  E=\bigwedge\nolimits^dV,
\]
选取 \(W\) 的一组基 \((w_1,\ldots,w_d)\)，再选取 \(E\) 的一组包含
\[
  e_0=w_1\wedge\cdots\wedge w_d.
\]
的基 \((e_0,\ldots,e_n)\)。\(G\) 在 \(V\) 上的作用典范地诱导 \(G\) 在 \(E\) 上的作用，因而诱导一个 \(A\)-余模结构
\[
  \rho:E\longrightarrow E\otimes A.
\]
对 \(j=0,\ldots,n\)，写成%
\begingroup
\renewcommand{\thefootnote}{\fnsymbol{footnote}}%
\footnote[2]{\label{note:sga3-VIB-11.16-rho-index-source-reading}
源文读法注：法文重版本和 Reinhold 对照版本均将求和下界印作 \(i=1\)，但所选基从 \(e_0\) 开始，且随后恒等式从 \(0\) 起求和。因此这里下界为 \(i=0\)。}
\endgroup
\phantomsection\label{form:sga3-VIB-11.16-rho-matrix}%
\[
  \rho(e_j)=\sum_{i=0}^{n}e_i\otimes b_{ij}.
\]
如 \hyperref[lem:sga3-VIB-11.12]{11.12} 的证明中一样，有
\phantomsection\label{form:sga3-VIB-11.16-star}%
\[
  \tag{\(*\)}
  \Delta(b_{ij})
    =\sum_{\ell=0}^{n}b_{i\ell}\otimes b_{\ell j}.
\]
置 \(a_i=b_{i0}\)。等价地，若 \((e_0^*,\ldots,e_n^*)\) 为对偶基，则各 \(a_i\) 为矩阵系数
\[
  g\longmapsto e_i^*(ge_0).
\]
\(H\) 在 \(E\) 上的作用对应于
\phantomsection\label{form:sga3-VIB-11.16-H-coaction-E}%
\[
  \overline{\rho}:\quad
  E\xrightarrow{\ \rho\ }E\otimes A
    \xrightarrow{\ \mathrm{id}_E\otimes\pi\ }E\otimes B.
\]
由于 \(W\) 在 \(H\) 下稳定，向量 \(e_0\) 关于 \(H\) 半不变。因此
\[
  \overline{\rho}(e_0)=e_0\otimes\pi(a_0),
  \qquad
  a_i=b_{i0}\in I\quad(1\leq i\leq n).
\]
代入 \hyperref[form:sga3-VIB-11.16-star]{\textup{(\(*\))}} 得到
\[
  \overline{\Delta}(a_i)
    =a_i\otimes\pi(a_0)
  \qquad(0\leq i\leq n).
\]
所以各 \(a_i\) 关于 \(H\) 半不变，其权均为 \(\pi(a_0)\)。此外，必要时以 \(e_0\) 生成的 \(G\)-子模替代 \(E\)\textup{（参见 \hyperref[lem:sga3-VIB-11.8]{11.8}）}，便可假设各 \(a_i\) 线性无关。

反之，设
\[
  H'=\operatorname{Spec}B',
  \qquad
  B'=A/I',
\]
是 \(G\) 的一个闭子群概形，使得每个 \(a_i\) 关于它半不变，其权为 \(\lambda_i\in B'\)。这尤其包括 \(e_0\) 关于 \(H'\) 半不变且权为 \(\lambda'\) 的情形。我们将证明 \(H'\subset H\)。%
\begingroup
\renewcommand{\thefootnote}{\fnsymbol{footnote}}%
\footnote[3]{\label{note:sga3-VIB-11.16-maximality-source-reading}
源文读法注：法文重版本此处写“证明 \(H'=H\)”，Reinhold 对照版本亦然。最大性断言所需的是 \(H'\subset H\)，而这正是下面所证明的结论。}
\endgroup\ 
设 \(\pi':A\to B'\) 为投影。由假设得到
\phantomsection\label{form:sga3-VIB-11.16-converse-coefficients}%
\[
  a_i\otimes\lambda_i
  =
  (\mathrm{id}_A\otimes\pi')\Delta(b_{i0})
  =
  \sum_{\ell=0}^{n}b_{i\ell}\otimes\pi'(a_\ell).
\]
可知
\[
  \lambda_i=\pi'(a_0),
  \qquad
  a_\ell\in I'\quad(1\leq\ell\leq n),
\]
从而 \(e_0\) 关于 \(H'\) 半不变。由\hyperref[lem:sga3-VIB-11.16.0]{引理 11.16.0}，\(W\) 在 \(H'\) 下稳定。由于 \(W\) 生成 \(I\)，理想 \(I\) 也在 \(H'\) 下稳定，因此
\[
  \Delta(I)\subset I\otimes A+A\otimes I'.
\]
最后，
\[
  I\subset\operatorname{Ker}\varepsilon,
  \qquad
  (\varepsilon\otimes\mathrm{id}_A)\circ\Delta
    =\mathrm{id}_A,
\]
故 \(I\subset I'\)，从而 \(H'\subset H\)。证明完毕。
